\documentclass[11pt,a4paper,UTF8]{ctexart}
\usepackage{amsmath,amssymb,mathtools}
\usepackage[dvipsnames]{xcolor}
\usepackage[most]{tcolorbox}
\usepackage{geometry}
\usepackage{booktabs,longtable,array,tabularx}
\usepackage{enumitem}
\usepackage{needspace}
\usepackage[colorlinks=true,linkcolor=MidnightBlue,urlcolor=MidnightBlue]{hyperref}

\geometry{left=2.1cm,right=2.1cm,top=2.2cm,bottom=2.2cm}
\setlength{\parindent}{0em}
\setlength{\parskip}{0.35em}
\linespread{1.08}
\setlist{nosep,leftmargin=2em}
\allowdisplaybreaks
\raggedbottom

\definecolor{bluebg}{RGB}{235,245,255}
\definecolor{bluebd}{RGB}{41,128,185}
\definecolor{greenbg}{RGB}{235,255,240}
\definecolor{greenbd}{RGB}{39,174,96}
\definecolor{orangebg}{RGB}{255,246,235}
\definecolor{orangebd}{RGB}{230,126,34}
\definecolor{redbg}{RGB}{255,240,240}
\definecolor{redbd}{RGB}{192,57,43}

\newtcolorbox{identifybox}[1][]{
  enhanced,breakable,colback=bluebg,colframe=bluebd,
  title=#1,fonttitle=\bfseries,arc=2mm,boxrule=0.8pt
}
\newtcolorbox{templatebox}[1][]{
  enhanced,breakable,colback=orangebg,colframe=orangebd,
  title=#1,fonttitle=\bfseries,arc=2mm,boxrule=0.9pt
}
\newtcolorbox{answerbox}[1][]{
  enhanced,breakable,colback=greenbg,colframe=greenbd,
  title=#1,fonttitle=\bfseries,arc=2mm,boxrule=0.8pt
}
\newtcolorbox{scorebox}[1][]{
  enhanced,breakable,colback=redbg,colframe=redbd,
  title=#1,fonttitle=\bfseries,arc=2mm,boxrule=0.8pt
}

\newcommand{\rank}{\operatorname{rank}}
\newcommand{\tr}{\operatorname{tr}}
\newcommand{\diag}{\operatorname{diag}}
\newcommand{\nul}{\operatorname{nul}}
\newcommand{\R}{\mathbb R}

\title{\Huge 四川大学线性代数（理工）期末真题融合讲义\\[0.25em]
\Large 2020--2025 历年真题按题型考点分类详解}
\author{按上传试卷与参考整理}
\date{\today}

\begin{document}
\maketitle
\thispagestyle{empty}
\tableofcontents
\newpage

\begin{scorebox}[定位]
版本：2020--2025 春季期末真题融合版\\
用法：每个板块按“识别题型 -> 套模板 -> 写计算 -> 得结论”组织；历年题目不放附录，而是嵌入对应考点。\\
\end{scorebox}

\bigskip\hrule\bigskip

\section{全题覆盖索引}

本稿覆盖你上传的 2020--2021、2021--2022、2022--2023、2023--2024、2024--2025 春季试卷全部题型，并参考 2019--2020 秋卷作补充。

{\scriptsize
\begin{longtable}{>{\raggedright\arraybackslash}p{0.13\textwidth}>{\raggedright\arraybackslash}p{0.25\textwidth}>{\raggedright\arraybackslash}p{0.43\textwidth}>{\raggedright\arraybackslash}p{0.13\textwidth}}
\toprule
年份 & 填空题 & 解答题 & 证明题 \\
\midrule
2020--2021 春 & 1--6：行列式、秩、余子式、矩阵幂、实对称特征、二次型秩 & 二：相似与矩阵表示；三：线性方程组；四：子空间与基；五：实对称矩阵；六：矩阵方程；七：二次型正定 & 八：基础解系变换；伴随矩阵与正交矩阵 \\
2021--2022 春 & 1--6：行列式列变换、矩阵幂、秩最小、特征值与伴随、坐标、正定条件 & 一：矩阵方程；二：方程组同解；三：极大无关组；四：转移矩阵模型；五：二次型正交化 & 一：可逆性；二：幂等矩阵相似标准形 \\
2022--2023 春 & 1--6：余子式、Sherman 型逆、矩阵幂、非齐次解维数、特征值、负定参数 & 二：向量组；三：矩阵方程；四：对角化；五：方程组；六：基与过渡矩阵；七：实对称二次型 & 一：伴随矩阵基础解系；二：二次多项式矩阵可对角化 \\
2023--2024 春 & 1--6：行列式常数项、合同、无穷多解、正交方程、坐标内积、特征值 & 二：伴随矩阵方程；三：向量组；四：方程组；五：一般基变换；六：实对称矩阵；七：二次型合同/正交 & 一：非齐次解组线性无关；二：迹与转置证明 \\
2024--2025 春 & 1--6：行列式、伴随、实对称惯性指数、过渡矩阵、解空间与秩、矩阵幂特征值 & 二：矩阵方程与矩阵幂；三：向量组；四：非齐次通解；五：基坐标；六：实对称矩阵；七：二次型 & 一：Vandermonde 型特征向量证明；二：秩一矩阵逆公式 \\
\bottomrule
\end{longtable}
}

\bigskip\hrule\bigskip

\section{行列式、伴随矩阵、秩：填空题先拿分}

\begin{identifybox}[高亮公式]
$|kA|=k^n|A|,\quad |A^*|=|A|^{n-1},\quad AA^*=A^*A=|A|E$\\
$r(AB)\le \min\{r(A),r(B)\}$；若 $P,Q$ 可逆，则 $r(PAQ)=r(A)$。\\
伴随常用：若 $A$ 可逆，则 $A^*=|A|A^{-1}$。\\
\end{identifybox}

\begin{templatebox}[考场模板]

\begin{enumerate}
\item 看是否是“可逆矩阵左右乘”：秩不变。
\item 看是否是“列向量线性组合”：把新列写成旧列乘一个系数矩阵，行列式相乘。
\item 看是否出现 $A^*$：立刻写 $A^*=|A|A^{-1}$ 或 $|A^*|=|A|^{n-1}$。
\item 行列式含参数或多项式：只求指定系数时，用列/行结构，不要硬展开整式。
\end{enumerate}

\end{templatebox}
\subsection{真题融合}

\begin{answerbox}[2020--2021 填空 1]
题：$A$ 为四阶矩阵，$|A|=\frac14$，求 $|(2A^*)^{-1}|$。

\textbf{【考点】} 基础计算

\textbf{【速解】}

四阶时 $|A^*|=|A|^3=\frac1{64}$，$\Rightarrow$ $|2A^*|=2^4\cdot\frac1{64}=\frac14$。
\textbf{【答案】}

\[|(2A^*)^{-1}|=4.\]
\end{answerbox}
\begin{answerbox}[2020--2021 填空 2]
题：设 $A$ 是 $4\times3$ 阶矩阵，$A$ 的秩 $r(A)=2$，
\[
B=\begin{pmatrix}
1&0&3\\
0&2&0\\
0&0&3
\end{pmatrix},
\]
则 $AB$ 的秩 $r(AB)=\underline{\qquad}$。

\textbf{【考点】} 基础计算

\textbf{【速解】}

先用 $r(AB)\le 2$，$B$ 的两个主元列经 $A$ 作用后仍线性无关，
\textbf{【答案】}

\[r(AB)=2.\]
\end{answerbox}
\begin{answerbox}[2020--2021 填空 3]
题：三阶矩阵 $A$ 第一行为 $(1,2,3)$，$|A|$ 第二行元素余子式分别为 $a+2,a+1,a-2$，求 $a$。

\textbf{【考点】} 基础计算

\textbf{【速解】}

按“某一行元素乘另一行余子式之和为 $0$”：
\textbf{【答案】}

\[1(a+2)+2(a+1)+3(a-2)=0,\]
得 $a=3$。
\end{answerbox}
\begin{answerbox}[2021--2022 填空 1]
题：设 $\alpha_1,\alpha_2,\alpha_3$ 均为 $3$ 维列向量，记矩阵
\[
A=(\alpha_1,\alpha_2,\alpha_3),
\]
矩阵
\[
B=(\alpha_1+\alpha_2+\alpha_3,\ \alpha_1+2\alpha_2+4\alpha_3,\ \alpha_1+3\alpha_2+9\alpha_3).
\]
如果 $|A|=1$，那么 $|B|=\underline{\qquad}$。

\textbf{【考点】} 基础计算

\textbf{【速解】}

写成 $B=AC$，只算系数矩阵 $C$ 的行列式；
\textbf{【答案】}

\[|B|=2.\]
\end{answerbox}
\begin{answerbox}[2021--2022 填空 3]
题：设矩阵
\[
A=\begin{pmatrix}
1&\lambda&-1&2\\
2&-1&\lambda&5\\
1&10&-6&1
\end{pmatrix},
\]
当参数 $\lambda=\underline{\qquad}$ 时，矩阵 $A$ 的秩最小。

\textbf{【考点】} 基础计算

\textbf{【速解】}

先令所有最高阶非零子式同时为 $0$，检查低阶子式不全为 $0$。
\textbf{【答案】}

\[\lambda=3.\]
\end{answerbox}
\begin{answerbox}[2021--2022 填空 4]
题：三阶矩阵 $A$ 的特征值为 $1,1,2$，求
\[
\left|\left(\frac13A^2\right)^{-1}+\frac12A^*-E\right|.
\]

\textbf{【考点】} 基础计算

\textbf{【速解】}

用 $A^*=|A|A^{-1}$，把问题转成特征值代入。
\textbf{【答案】}

\[\frac94.\]
\end{answerbox}
\begin{answerbox}[2022--2023 填空 1]
题：设
\[
A=\begin{pmatrix}
1&2&4&2\\
1&3&9&0\\
1&-1&1&2\\
1&4&16&3
\end{pmatrix},
\]
$M_{ij}$ 为矩阵 $A$ 的第 $i$ 行第 $j$ 列元素的余子式，$i,j=1,2,3,4$，则
\[
-8M_{14}+27M_{24}+M_{34}+64M_{44}=\underline{\qquad}.
\]

\textbf{【考点】} 基础计算

\textbf{【速解】}

这是按第 4 列余子式反造行列式。把第 4 列替换成 $(8,27,-1,64)^T$ 后按第 4 列展开，用列运算求值。
\textbf{【答案】}

\[120.\]
\end{answerbox}
\begin{answerbox}[2022--2023 填空 2]
题：$\alpha=(c,0,\cdots,0,c)^T,c>0$，$A=E-\alpha\alpha^T$，$A^{-1}=E+c^{-1}\alpha\alpha^T$，求 $c$。

\textbf{【考点】} 基础计算

\textbf{【速解】}

套
\[
(E-\alpha\alpha^T)^{-1}
=E+\frac{\alpha\alpha^T}{1-\alpha^T\alpha}.
\]
而 $\alpha^T\alpha=2c^2$，与 $E+c^{-1}\alpha\alpha^T$ 比较：
\textbf{【答案】}

\[
\frac1{1-2c^2}=\frac1c
\Rightarrow c=1-2c^2
\Rightarrow c=\frac12.
\]
\end{answerbox}
\begin{answerbox}[2023--2024 填空 1]
题：多项式
\[
f(x)=
\begin{vmatrix}
x&x&1&2x\\
1&x&2&-1\\
2&1&x&1\\
2&-1&1&x
\end{vmatrix}
\]
的常数项。

\textbf{【考点】} 基础计算

\textbf{【速解】}

常数项就是令 $x=0$ 后的行列式：
\textbf{【答案】}

\[
\begin{vmatrix}
0&0&1&0\\
1&0&2&-1\\
2&1&0&1\\
2&-1&1&0
\end{vmatrix}=5.
\]
\end{answerbox}
\begin{answerbox}[2023--2024 填空 2]
题：若
\[
A=\begin{pmatrix}a&-1&-1\\-1&a&-1\\-1&-1&a\end{pmatrix}
\]
与
\[
B=\begin{pmatrix}1&1&0\\0&-1&1\\1&0&1\end{pmatrix}
\]
等价，即 $A$ 可经有限次初等变换化为 $B$，则 $a=\underline{\qquad}$。

\textbf{【考点】} 基础计算

\textbf{【速解】}

等价矩阵秩相同。由标准，先化简 $B$ 得 $r(B)=2$；令 $r(A)=2$，即 $|A|=0$ 且二阶子式不全为 $0$，得
\textbf{【答案】}

\[a=2.\]
\end{answerbox}
\begin{answerbox}[2024--2025 填空 1]
题：计算四阶行列式
\[
\begin{vmatrix}
3&1&1&2\\
-1&2&2&2\\
1&2&-1&1\\
-3&2&0&1
\end{vmatrix}.
\]

\textbf{【考点】} 基础计算

\textbf{【速解】}

用行变换化上三角，注意交换行变号、倍加不变。结果：
\textbf{【答案】}

\[-13.\]
\end{answerbox}
\begin{answerbox}[2024--2025 填空 2]
题：三阶方阵 $|A|=\frac13$，求 $|A^*|+|(2A)^{-1}-A^*|$。

\textbf{【考点】} 基础计算

\textbf{【速解】}

\[|A^*|=|A|^2=\frac19,\quad A^*=\frac13 A^{-1},\quad (2A)^{-1}=\frac12A^{-1}.\]
于是
\[|(2A)^{-1}-A^*|=\left|\frac16A^{-1}\right|=\frac1{6^3}\cdot 3=\frac1{72},\]

\textbf{【答案】}

\[\frac19+\frac1{72}=\frac18.\]
\end{answerbox}
\bigskip\hrule\bigskip

\section{矩阵方程与矩阵幂}

\begin{identifybox}[高亮公式]
$XA+M=N+2X \Rightarrow X(A-2E)=N-M$\\
$A=PBP^{-1}\Rightarrow A^n=PB^nP^{-1}$\\
若 $A^2=A$，则 $A^n=A\ (n\ge1)$。\\
\end{identifybox}

\begin{templatebox}[考场模板]

\begin{enumerate}
\item 把未知矩阵 $X$ 全部移到一边。
\item 提出右乘或左乘的公共矩阵。
\item 求逆时优先用已经给出的三角矩阵、伴随矩阵或相似关系。
\item 求高次幂时先判断：幂等、对角矩阵、相似对角化、递推关系。
\end{enumerate}

\end{templatebox}
\subsection{真题融合}

\begin{answerbox}[2020--2021 解答二：相似表示与 |A+E|]
题：$P=(\alpha,A\alpha,A^2\alpha)$ 可逆，且 $A^3\alpha=5A\alpha-4A^2\alpha$。求 $B$ 使 $A=PBP^{-1}$，求 $|A+E|$。

\textbf{【题目分析】}

\[
AP=(A\alpha,A^2\alpha,A^3\alpha)=(\alpha,A\alpha,A^2\alpha)
\begin{pmatrix}
0&0&0\\
1&0&5\\
0&1&-4
\end{pmatrix}.
\]

\textbf{【详细解答】}

化简得
\[B=\begin{pmatrix}0&0&0\\1&0&5\\0&1&-4\end{pmatrix}.\]

\textbf{【最终答案】}

由 $A$ 与 $B$ 相似：
\[|A+E|=|B+E|=-8.\]
\end{answerbox}
\begin{answerbox}[2020--2021 解答六：矩阵方程 XA+2B=AB+2X]
题：设
\[
A=\begin{pmatrix}3&0&0\\2&1&0\\2&1&3\end{pmatrix},\quad
B=\begin{pmatrix}1&0&0\\0&1&0\\0&0&-1\end{pmatrix},
\]
矩阵 $X$ 满足
\[
XA+2B=AB+2X.
\]
（1）求 $(A-2E)^{-1}$；（2）求 $X^{2021}$。

\textbf{【题目分析】}

矩阵方程中 $X$ 在左侧，应整理成右乘公共矩阵后再右乘逆矩阵。

\textbf{【详细解答】}

参考答案行化简：
\[
(A-2E)^{-1}=
\begin{pmatrix}
-1&2&-4\\
0&1&-1\\
0&0&1
\end{pmatrix}.
\]
由
\[
XA+2B=AB+2X
\]
得
\[
X(A-2E)=(A-2E)B.
\]
右乘 \((A-2E)^{-1}\)，得
\[
X=(A-2E)B(A-2E)^{-1}.
\]
因此
\[
X^{2021}=(A-2E)B^{2021}(A-2E)^{-1}.
\]

\textbf{【最终答案】}

\[
X^{2021}=
\begin{pmatrix}
1&0&-8\\
0&1&-2\\
0&0&1
\end{pmatrix}.
\]

\textbf{【方法总结】}

矩阵方程中若得到 \(X(A-2E)=M\)，只能右乘 \((A-2E)^{-1}\)，不能左乘或交换顺序。
\end{answerbox}
\begin{answerbox}[2020--2021 填空 4：置换矩阵幂]
题：
\[
P=\begin{pmatrix}0&1&0\\1&0&0\\0&0&1\end{pmatrix},\quad
Q=\begin{pmatrix}2&5&12\\1&4&9\\3&6&15\end{pmatrix},\quad
R=\begin{pmatrix}1&0&-2\\0&1&0\\0&0&1\end{pmatrix},
\]
求 $P^{2021}QR$。

\textbf{【考点】} 矩阵幂

\textbf{【速解】}

$P^2=E$，$\Rightarrow$ $P^{2021}=P$。
\textbf{【答案】}

\[
P^{2021}QR=
\begin{pmatrix}
1&4&7\\
2&5&8\\
3&6&9
\end{pmatrix}.
\]
\end{answerbox}
\begin{answerbox}[2021--2022 解答一：低秩矩阵方程]
题：设向量 $\alpha=(1,2)^T$，$\beta=\left(1,\frac12\right)^T$，
\[
\gamma=\begin{pmatrix}8&4\\16&8\end{pmatrix},
\]
令 $A=\alpha\beta^T,\ B=\beta^T\alpha$，其中 $\beta^T$ 表示列向量 $\beta$ 的转置，求解方程
\[
2B^2A^2X=A^4X+B^3X+\gamma.
\]

\textbf{【题目分析】}

计算
\[A^2=(\alpha^T\alpha)A,\]

\textbf{【详细解答】}

计算
\[A^2=(\alpha^T\alpha)A,\]

\textbf{【最终答案】}


\[
X=
\begin{pmatrix}
1&\frac12\\
2&1
\end{pmatrix}.
\]

\textbf{【方法总结】}

低秩矩阵题先用 \(A^2=(\alpha^T\alpha)A\) 降幂，再代回原方程求 \(X\)。
\end{answerbox}
\begin{answerbox}[2021--2022 填空 2]
题：设
\[
A=\begin{pmatrix}
1&0&1\\
0&2&0\\
1&0&1
\end{pmatrix},
\]
则 $A^{2022}-2A^{2021}=\underline{\qquad}$。

\textbf{【考点】} 基础计算

\textbf{【速解】}

先验证 $A^2=2A$，$\Rightarrow$ $A^{2022}=2A^{2021}$，
\textbf{【答案】}

\[O.\]
\end{answerbox}
\begin{answerbox}[2022--2023 填空 3]
题：已知三阶方阵
\[
A=\begin{pmatrix}0&1&0\\1&0&0\\0&0&1\end{pmatrix},\quad
B=\begin{pmatrix}1&0&0\\0&1&2\\0&0&1\end{pmatrix},
\]
则矩阵 $A^{2023}B^{2024}$ 位于第二行第三列的元素的值为 $\underline{\qquad}$。

\textbf{【考点】} 基础计算

\textbf{【速解】}


\[
A=\begin{pmatrix}0&1&0\\1&0&0\\0&0&1\end{pmatrix},\quad
B=\begin{pmatrix}1&0&0\\0&1&2\\0&0&1\end{pmatrix}.
\]
$A^2=E$，$\Rightarrow$ $A^{2023}=A$；$B=E+N,\ N^2=0$，$\Rightarrow$ $B^{2024}=E+2024N$。乘后取第二行第三列，
\textbf{【答案】}

\[0.\]
\end{answerbox}
\begin{answerbox}[2022--2023 解答三：矩阵方程]
题：已知方阵
\[
A=\begin{pmatrix}
1&1&0&0\\
1&2&0&0\\
0&0&0&1\\
0&0&2&0
\end{pmatrix}
\]
满足
\[
\left[(|A-E|^{-1}A)^*\right]^{-1}XA^{-1}=E-\frac12AX,
\]
求矩阵 $X$。

\textbf{【题目分析】}

计算
\[|A|=-2,\quad |A-E|=1.\]

\textbf{【详细解答】}

原式化为
\[X=2(A-E)^{-1}.\]

\textbf{【最终答案】}


\[
X=
\begin{pmatrix}
-2&0&0&0\\
2&0&0&0\\
0&0&2&2\\
0&0&4&2
\end{pmatrix}.
\]
\end{answerbox}
\begin{answerbox}[2021--2022 解答四：转移矩阵模型]
题：某城市及郊区乡镇共有 $50$ 万人从事农、工、商工作，假设总人数在若干年内保持稳定不变。当前约有 $25$ 万人从事农业，约 $15$ 万人从事工业，约 $10$ 万人从事商业；在务农人员中，每年约有 $20\%$ 改为务工，$10\%$ 改为经商；在务工人员中，每年约有 $20\%$ 改为务农，$10\%$ 改为经商；在经商人员中，每年约有 $10\%$ 改为务农，$10\%$ 改为务工。
（1）请写出 $1$ 年后从事农、工、商工作的人数分别是多少；
（2）请预测 $n$ 年之后从事各行业人员总数及发展趋势。

\textbf{【题目分析】}

写状态转移矩阵
\[
A=\begin{pmatrix}
0.7&0.2&0.1\\
0.2&0.7&0.1\\
0.1&0.1&0.8
\end{pmatrix}.
\]

\textbf{【详细解答】}

一年后
\[
A(25,15,10)^T=(21.5,16.5,12)^T.
\]

\textbf{【最终答案】}

因 $1$ 是主特征值，其余特征值绝对值小于 $1$，长期趋于稳定分布：
\[\left(\frac{50}{3},\frac{50}{3},\frac{50}{3}\right)^T.\]
\end{answerbox}
\begin{answerbox}[2023--2024 解答二：伴随矩阵方程]
题：
\[
A=\begin{pmatrix}1&0&0\\0&1&0\\0&1&1\end{pmatrix},\quad
B=\begin{pmatrix}1&0&0\\0&0&1\\0&1&0\end{pmatrix},
\]
且 $A^*X+B=2BA+2X$，求 $X$。

\textbf{【题目分析】}

\[A^*X-2X=2BA-B,\quad (A^*-2E)X=2BA-B.\]

\textbf{【详细解答】}

因此
\[X=(A^*-2E)^{-1}(2BA-B).\]

\textbf{【最终答案】}

参考答案给出
\[
A^*=
\begin{pmatrix}
1&0&0\\
0&1&0\\
0&-1&1
\end{pmatrix},
\]
故
\[
X=
\begin{pmatrix}
-1&0&0\\
0&-2&-1\\
0&1&1
\end{pmatrix}.
\]

\textbf{【方法总结】}

这里 \(X\) 在右侧公共因子之后，整理为 \((A^*-2E)X=\cdots\)，所以左乘 \((A^*-2E)^{-1}\)。
\end{answerbox}
\begin{answerbox}[2024--2025 解答二：BA+CB=\ensuremath{B^2} 与 \ensuremath{A^{2025}}]
题：
\[
B=\begin{pmatrix}1&0&0\\1&1&0\\1&1&1\end{pmatrix},\quad
C=\begin{pmatrix}0&0&0\\1&0&0\\1&1&1\end{pmatrix},
\]
满足 $BA+CB=B^2$，求 $A,A^{2025}$。

\textbf{【题目分析】}

\[BA=B^2-CB,\quad A=B^{-1}(B^2-CB).\]

\textbf{【详细解答】}

计算得
\[
A=\begin{pmatrix}
1&0&0\\
0&1&0\\
-1&-1&0
\end{pmatrix}.
\]

\textbf{【最终答案】}

直接验证 $A^2=A$，得
\[A^{2025}=A.\]
\end{answerbox}
\bigskip\hrule\bigskip

\section{向量组、秩、极大线性无关组}

\begin{identifybox}[高亮公式]
列向量组问题统一写矩阵 $A=(\alpha_1,\ldots,\alpha_s)$。\\
行最简形主元列对应原矩阵主元列；零空间维数 $=n-r(A)$。\\
\end{identifybox}

\begin{templatebox}[考场模板]

\begin{enumerate}
\item 把向量按列排成矩阵。
\item 初等行变换到阶梯形。
\item 主元列给极大无关组。
\item 若问某向量能否由其余向量表示：看对应列是否非主元列，并从行最简形读系数。
\end{enumerate}

\end{templatebox}
\subsection{真题融合}

\begin{answerbox}[2021--2022 解答三：极大无关组与零空间]
题：设矩阵
\[
A=\begin{pmatrix}
1&-2&-1&1&-7\\
1&-2&2&3&-1\\
2&-4&1&8&-4
\end{pmatrix}.
\]
（1）请给出矩阵 $A$ 的列向量组的一个极大无关组；
（2）矩阵 $A$ 的零空间 $\operatorname{nul}(A)$ 的维数是多少？请给出零空间 $\operatorname{nul}(A)$ 的一组最小生成集。

\textbf{【题目分析】}

本题考查列向量组极大无关组与零空间。将列向量组成矩阵 \(A=(\alpha_1,\ldots,\alpha_5)\)，行化简后用主元列确定极大无关组。

\textbf{【详细解答】}

参考答案行最简形显示主元列为第 \(1,3,4\) 列，故
\[
r(A)=3.
\]
主元列对应原矩阵列向量，不是化简后矩阵的列向量。

\textbf{【最终答案】}

一个极大线性无关组为
\[
\alpha_1,\alpha_3,\alpha_4.
\]
零空间维数为
\[
\nul(A)=5-r(A)=2.
\]
零空间的一组最小生成集可取
\[
\xi_1=(2,1,0,0,0)^T,\qquad
\xi_2=\left(\frac{20}{3},0,-\frac43,-1,1\right)^T.
\]

\textbf{【方法总结】}

极大无关组必须由原矩阵的主元列给出；零空间维数等于列数减秩。
\end{answerbox}
\begin{answerbox}[2022--2023 解答二：含参数向量组]
题：已知向量组
\[
\alpha_1=(1,0,1,0)^T,\quad
\alpha_2=(3,-1,2,1)^T,\quad
\alpha_3=(1,a,b,0)^T,\quad
\alpha_4=(1,2,-2,1)^T.
\]
（1）求向量组的秩和一个极大线性无关组；
（2）$\alpha_3$ 能否由其余向量线性表示？若能，请给出线性表示的表达式。

\textbf{【题目分析】}

把四个向量作列，按化简：
\[
(\alpha_1,\alpha_2,\alpha_4,\alpha_3)
\sim
\begin{pmatrix}
1&3&1&1\\
0&1&1&0\\
0&0&3&a\\
0&0&0&2a+3b-3
\end{pmatrix}.
\]

\textbf{【详细解答】}

把四个向量作列，按化简：
\[
(\alpha_1,\alpha_2,\alpha_4,\alpha_3)
\sim
\begin{pmatrix}
1&3&1&1\\
0&1&1&0\\
0&0&3&a\\
0&0&0&2a+3b-3
\end{pmatrix}.
\]

\textbf{【最终答案】}

分类讨论：
\begin{enumerate}
\item 若 $2a+3b\ne3$，则 $r=4$，极大无关组为 $\alpha_1,\alpha_2,\alpha_3,\alpha_4$，$\alpha_3$ 不能由其余向量表示。
\item 若 $2a+3b=3$，则 $r=3$，取极大无关组 $\alpha_1,\alpha_2,\alpha_4$，且
\end{enumerate}
\[
\alpha_3=\left(1+\frac{2a}{3}\right)\alpha_1-\frac a3\alpha_2+\frac a3\alpha_4.
\]
\end{answerbox}
\begin{answerbox}[2023--2024 解答三：参数 p 与秩]
题：
\[
\alpha_1=(1,1,2,3)^T,\ \alpha_2=(-1,-3,0,5)^T,\
\alpha_3=(3,2,p+2,-1)^T,\ \alpha_4=(-2,-6,p,10)^T.
\]
问 $p$ 何值时相关，并求秩与全部极大无关组。

\textbf{【题目分析】}

构造 $A=(\alpha_1,\alpha_2,\alpha_3,\alpha_4)$，令 $|A|=0$ 得
\[p=0.\]

\textbf{【详细解答】}

构造 $A=(\alpha_1,\alpha_2,\alpha_3,\alpha_4)$，令 $|A|=0$ 得
\[p=0.\]

\textbf{【最终答案】}

代回行化简，$r(A)=3$，全部极大无关组为
\[\alpha_1,\alpha_2,\alpha_3;\qquad \alpha_1,\alpha_3,\alpha_4.\]
\end{answerbox}
\begin{answerbox}[2024--2025 解答三：参数 \ensuremath{\mu} 与全部极大无关组]
题：
\[
(1,2,0,3)^T,\ (-1,-4,6,1)^T,\ (3,5,-4,\mu+5)^T,\ (-2,-8,12,\mu+3)^T
\]
线性相关，求 $\mu$ 与全部极大无关组。

\textbf{【题目分析】}

设四向量为列，行列式为零：
\[|A|=0\Rightarrow \mu=-1.\]

\textbf{【详细解答】}

设四向量为列，行列式为零：
\[|A|=0\Rightarrow \mu=-1.\]

\textbf{【最终答案】}

代回后 $r(A)=3$，行最简形显示第 $4$ 列可由第 $2$ 列表示，极大无关组为
\[\alpha_1,\alpha_2,\alpha_3;\qquad \alpha_1,\alpha_3,\alpha_4.\]
\end{answerbox}
\begin{answerbox}[2020--2021 证明八(1)：基础解系变换]
题：设向量组 $\alpha_1,\alpha_2,\alpha_3$ 是齐次线性方程组 $AX=0$ 的基础解系，证明：
\[
\alpha_1+\alpha_2+\alpha_3,\quad
\alpha_1+9\alpha_2+16\alpha_3,\quad
\alpha_1-3\alpha_2-4\alpha_3
\]
也是 $AX=0$ 的基础解系。

\textbf{【证明思路】}

先说明新向量都是原解的线性组合，$\Rightarrow$仍是解；写过渡矩阵 $C$，证明
\[|C|\ne0.\]

\textbf{【详细证明】}

先说明新向量都是原解的线性组合，$\Rightarrow$仍是解；写过渡矩阵 $C$，证明
\[|C|\ne0.\]

\textbf{【最终答案】}

$\Rightarrow$新向量组线性无关，个数又等于零空间维数，$\Rightarrow$仍为基础解系。
\end{answerbox}
\Needspace{20\baselineskip}
\bigskip\hrule\bigskip

\section{线性方程组：有解、无穷多解、通解}

\begin{identifybox}[高亮公式]
$Ax=b$ 有解 $\Longleftrightarrow r(A)=r(A,b)$。\\
唯一解 $\Longleftrightarrow r(A)=r(A,b)=n$。\\
无穷多解 $\Longleftrightarrow r(A)=r(A,b)<n$。\\
通解：$x=\eta_0+k_1\xi_1+\cdots+k_{n-r}\xi_{n-r}$。\\
\end{identifybox}

\begin{templatebox}[考场模板]

\begin{enumerate}
\item 写增广矩阵 $(A,b)$。
\item 行化简，比较 $r(A)$ 与 $r(A,b)$。
\item 有参数时先分母/主元为零处分情况讨论。
\item 通解必须写成“一个特解 + 导出组基础解系”。
\end{enumerate}

\end{templatebox}
\subsection{真题融合}

\begin{answerbox}[2020--2021 解答三：含参数三元方程组]
题：线性方程组
\[
\begin{cases}
x_1+x_2+3x_3=5,\\
x_1+2x_2+4x_3=7,\\
x_1+(a+3)x_2+(a^2+3)x_3=3a+10,
\end{cases}
\]
其中 $a$ 为常数。
（1）写出该方程组的增广矩阵；
（2）$a$ 为何值时方程组有解？有解时求出所有的解。

\textbf{【题目分析】}

方程组为
\[
\begin{cases}
x_1+x_2+3x_3=5,\\
x_1+2x_2+4x_3=7,\\
x_1+(a+3)x_2+(a^2+3)x_3=3a+10.
\end{cases}
\]

\textbf{【详细解答】}

构造增广矩阵：
\[
(A|b)=
\begin{pmatrix}
1&1&3&5\\
1&2&4&7\\
1&a+3&a^2+3&3a+10
\end{pmatrix}.
\]
行化简得
\[
(A|b)\sim
\begin{pmatrix}
1&1&3&5\\
0&1&1&2\\
0&0&(a-2)(a+1)&a+1
\end{pmatrix}.
\]

分类讨论：
\begin{enumerate}
\item $a\ne2,-1$ 时，
\[
r(A)=r(A|b)=3,
\]
方程组有唯一解
\end{enumerate}
\[
x=\left(\frac{3a-8}{a-2},\frac{2a-5}{a-2},\frac1{a-2}\right)^T.
\]

\begin{enumerate}
\item $a=-1$ 时，
\[
(A|b)\sim
\begin{pmatrix}
1&0&2&3\\
0&1&1&2\\
0&0&0&0
\end{pmatrix},
\qquad r(A)=r(A|b)=2<3,
\]
方程组有无穷多解
\end{enumerate}
\[
x=(3,2,0)^T+k(-2,-1,1)^T,\quad k\in\mathbb R.
\]

\textbf{【最终答案】}

\begin{enumerate}
\item $a\ne2,-1$ 时，
\[
x=\left(\frac{3a-8}{a-2},\frac{2a-5}{a-2},\frac1{a-2}\right)^T.
\]
\item $a=-1$ 时，
\[
x=(3,2,0)^T+k(-2,-1,1)^T,\quad k\in\mathbb R.
\]
\item $a=2$ 时，
\[
r(A)=2,\qquad r(A|b)=3,
\]
方程组无解。
\end{enumerate}

\textbf{【方法总结】}

含参数线性方程组必须先写增广矩阵，再由主元 \((a-2)(a+1)\) 分情况讨论；无穷多解时通解写成“特解 + 齐次基础解系”。
\end{answerbox}
\begin{answerbox}[2021--2022 解答二：两个方程组同解]
题：考虑如下非齐次方程组
\[
I:\begin{cases}
x_1+x_2-2x_4=-6,\\
4x_1-x_2-x_3-x_4=1,\\
3x_1-x_2-x_3=3,
\end{cases}
\quad
II:\begin{cases}
x_1+ax_2-x_3-x_4=-5,\\
bx_2-x_3-2x_4=-11,\\
x_3-2x_4=-c+1.
\end{cases}
\]
（1）求解方程组 $I$，用导出组基础解系表示通解；
（2）方程组 $I$ 与方程组 $II$ 同解，求参数 $a,b,c$ 的值。

\textbf{【题目分析】}

第一步化简 I：
\[
x=(-2,-4,-5,0)^T+k(1,1,2,1)^T.
\]

\textbf{【详细解答】}

第一步化简 I：
\[
x=(-2,-4,-5,0)^T+k(1,1,2,1)^T.
\]

\textbf{【最终答案】}

第二步：同解的意思是 II 必须对上面通解恒成立。把通解代入 II，比较常数项与 $k$ 的系数：
\[a=2,\quad b=4,\quad c=6.\]
\end{answerbox}
\begin{answerbox}[2022--2023 填空 4]
题：$n$ 阶方阵 $A$ 作系数矩阵，非齐次方程组 $Ax=b$ 仅有两个线性无关解，求 $r(A^*)$。

\textbf{【考点】} 基础计算

\textbf{【速解】}

非齐次解集不是线性空间；“仅有两个线性无关解”说明导出组解空间维数为 $1$，得
\[n-r(A)=1,\quad r(A)=n-1.\]
故 $A$ 不可逆且秩为 $n-1$。
\textbf{【答案】}

\[r(A^*)=1.\]
\end{answerbox}
\begin{answerbox}[2022--2023 解答五：非齐次方程组有三个无关解]
题：已知非齐次线性方程组
\[
\begin{cases}
x_1+x_2+x_3+x_4=-1,\\
4x_1+3x_2+5x_3-x_4=-1,\\
ax_1+x_2+3x_3+bx_4=1
\end{cases}
\]
有三个线性无关的解。
（1）求 $a,b$ 的值；
（2）求线性方程组的通解。

\textbf{【题目分析】}

三个线性无关解推出非齐次解集方向空间维数为 \(2\)，故 \(n-r(A)=2\)，于是 \(r(A)=2\)。

\textbf{【详细解答】}

参考答案对系数矩阵行化简：
\[
\begin{pmatrix}
1&1&1&1\\
4&3&5&-1\\
a&1&3&b
\end{pmatrix}
\sim
\begin{pmatrix}
1&1&1&1\\
0&2&1&-4\\
0&0&a-2&b+3
\end{pmatrix}.
\]
由 \(r(A)=2\) 得
\[
a=2,\qquad b=-3.
\]
此时增广矩阵化为
\[
(A|b)\sim
\begin{pmatrix}
1&0&2&-4&2\\
0&1&-1&5&-3\\
0&0&0&0&0
\end{pmatrix}.
\]

\textbf{【最终答案】}

主元变量为 \(x_1,x_2\)，自由变量为 \(x_3,x_4\)。通解为
\[
X=
\begin{pmatrix}2\\-3\\0\\0\end{pmatrix}
k_1\begin{pmatrix}-2\\1\\1\\0\end{pmatrix}
k_2\begin{pmatrix}4\\-5\\0\\1\end{pmatrix},
\qquad k_1,k_2\in\mathbb R.
\]

\textbf{【方法总结】}

非齐次通解必须写成“特解 + 齐次方程组基础解系”，不能只写模板。
\end{answerbox}
\begin{answerbox}[2023--2024 填空 3]
题：
\[A=\begin{pmatrix}1&0&-1\\1&1&-1\\0&1&a^2-1\end{pmatrix},\quad b=(0,1,a)^T,\]
若 $Ax=b$ 有无穷多解，求 $a$。

\textbf{【考点】} 基础计算

\textbf{【速解】}

无穷多解要求
\[|A|=0,\quad r(A)=r(A,b)<3.\]
代入检查可得
\textbf{【答案】}

\[a=-1.\]
\end{answerbox}
\begin{answerbox}[2023--2024 填空 4]
题：$A=(a_{ij})$ 是三阶正交矩阵，$a_{11}=1$，$b=(1,0,0)^T$，求 $Ax=b$ 的解。

\textbf{【考点】} 基础计算

\textbf{【速解】}

正交矩阵满足 $A^{-1}=A^T$。由 $a_{11}=1$ 且第一行长度为 $1$，第一行只能是 $(1,0,0)$；同理第一列也是 $(1,0,0)^T$。因此
\textbf{【答案】}

\[
x=A^Tb=(1,0,0)^T.
\]
\end{answerbox}
\begin{answerbox}[2023--2024 解答四]
题：
\[
A=\begin{pmatrix}\lambda&1&1\\0&\lambda-1&0\\1&1&\lambda\end{pmatrix},\quad
b=(a,1,1)^T,
\]
且 $Ax=b$ 有两个不同解。求 $\lambda,a$ 与通解。

\textbf{【题目分析】}

两个不同解即无穷多解：
\[|A|=0,\quad r(A)=r(A,b)<3.\]

\textbf{【详细解答】}

由增广矩阵化简得：
\[\lambda=-1,\quad a=-2.\]

\textbf{【最终答案】}

此时
\[
(A|b)\sim
\begin{pmatrix}
1&0&-1&\frac32\\
0&1&0&-\frac12\\
0&0&0&0
\end{pmatrix}.
\]
主元变量为 \(x_1,x_2\)，自由变量为 \(x_3\)。通解为
\[
x=
\begin{pmatrix}
\frac32\\[0.2em]
-\frac12\\[0.2em]
0
\end{pmatrix}
k
\begin{pmatrix}
1\\0\\1
\end{pmatrix},
\qquad k\in\mathbb R.
\]

\textbf{【方法总结】}

有两个不同解说明无穷多解，必须同时满足 \(r(A)=r(A|b)<n\)。
\end{answerbox}
\begin{answerbox}[2024--2025 填空 5]
题：非零三阶方阵 $A$ 的每一列都是齐次方程组
\[
\begin{cases}
x_1+2x_2-2x_3=0,\\
3x_1+x_2+\lambda x_3=0,\\
4x_1+3x_2-3x_3=0
\end{cases}
\]
的解，求 $\lambda+r(A)$。

\textbf{【考点】} 基础计算

\textbf{【速解】}

$A\ne O$ 且每列在解空间中，$\Rightarrow$齐次方程组有非零解：
\[
\left|\begin{matrix}1&2&-2\\3&1&\lambda\\4&3&-3\end{matrix}\right|=0
\Rightarrow \lambda=-1.
\]
此时系数矩阵秩为 $2$，解空间维数为 $1$，$\Rightarrow$非零矩阵 $A$ 的列均在一维空间中，$r(A)=1$。
\textbf{【答案】}

\[\lambda+r(A)=0.\]
\end{answerbox}
\begin{answerbox}[2024--2025 解答四：由三个解求通解]
题：$r(A)=2$，$\eta_1,\eta_2,\eta_3$ 是 $Ax=b$ 的三个解，且
\[4\eta_1-\eta_2=(3,0,-3)^T,\quad 5\eta_1-3\eta_3=(-4,12,2)^T.\]
求通解。

\textbf{【题目分析】}

三阶方阵 \(A\) 满足 \(r(A)=2\)，故齐次方程组 \(Ax=0\) 的基础解系由
\[
n-r(A)=3-2=1
\]
个解向量构成。非齐次方程组解的仿射组合仍为非齐次方程组的解。

\textbf{【详细解答】}

由题设
\[
4\eta_1-\eta_2=
\begin{pmatrix}3\\0\\-3\end{pmatrix},
\qquad
5\eta_1-3\eta_3=
\begin{pmatrix}-4\\12\\2\end{pmatrix}.
\]
因为 \(4-1=3\)、\(5-3=2\)，所以
\[
\frac13(4\eta_1-\eta_2)=
\begin{pmatrix}1\\0\\-1\end{pmatrix},
\qquad
\frac12(5\eta_1-3\eta_3)=
\begin{pmatrix}-2\\6\\1\end{pmatrix}
\]
都是 \(Ax=b\) 的解。

两解相减得
\[
\begin{pmatrix}3\\-6\\-2\end{pmatrix}
=
\begin{pmatrix}1\\0\\-1\end{pmatrix}
-
\begin{pmatrix}-2\\6\\1\end{pmatrix},
\]
这是 \(Ax=0\) 的一个非零解。又 \(Ax=0\) 的基础解系只含一个向量，故它可作为基础解系。

\textbf{【最终答案】}

取特解
\[
\eta_0=
\begin{pmatrix}1\\0\\-1\end{pmatrix},
\qquad
\xi=
\begin{pmatrix}3\\-6\\-2\end{pmatrix}.
\]
所以 \(Ax=b\) 的通解为
\[
x=
k\begin{pmatrix}3\\-6\\-2\end{pmatrix}
+
\begin{pmatrix}1\\0\\-1\end{pmatrix},
\qquad k\in\mathbb R.
\]

\textbf{【方法总结】}

本题按参考答案写法：先构造两个非齐次特解，再用“两个特解之差”得到齐次方程组的基础解系。
\end{answerbox}
\bigskip\hrule\bigskip

\section{基、坐标、过渡矩阵}

\begin{identifybox}[高亮公式]
若 $(\beta_1,\ldots,\beta_n)=(\alpha_1,\ldots,\alpha_n)C$，则 $C$ 是“新基向量在旧基下的坐标矩阵”。\\
同一向量：$[v]_\alpha=C[v]_\beta$，$\Rightarrow$ $[v]_\beta=C^{-1}[v]_\alpha$。\\
\end{identifybox}

\begin{templatebox}[考场模板]

\begin{enumerate}
\item 先写 $(\beta)=(\alpha)C$。
\item 判断是否为基：只看 $|C|\ne0$。
\item 坐标变换：旧坐标 $=C$ 新坐标；新坐标 $=C^{-1}$ 旧坐标。
\item “两组基下坐标相同”：设共同坐标为 $x$，写 $(\alpha)x=(\beta)x$，化成 $(C-E)x=0$。
\end{enumerate}

\end{templatebox}
\subsection{真题融合}

\begin{answerbox}[2020--2021 解答四：子空间、基与过渡矩阵]
题：设 $V$ 是由
\[
\alpha_1=(2,3,1,1)^T,\quad \alpha_2=(1,2,-1,1)^T
\]
所张成的子空间，
\[
\beta_1=(4,7,-1,3)^T,\quad \beta_2=(3,4,3,1)^T.
\]
（1）求子空间 $V$ 的维数 $\dim(V)$；
（2）判断是否有 $\beta_1\in V,\ \beta_2\in V$，并说明理由；
（3）$\beta_1,\beta_2$ 是否是子空间 $V$ 的基？请说明理由。如果是，求基 $\beta_1,\beta_2$ 到基 $\alpha_1,\alpha_2$ 的过渡矩阵。

\textbf{【题目分析】}

把生成向量作列行化简，卷得到
\[\dim V=2.\]

\textbf{【详细解答】}

把待判向量并入列，若秩不变则属于 $V$。由
\[(\beta_1,\beta_2)=(\alpha_1,\alpha_2)C\]

\textbf{【最终答案】}

直接读出过渡矩阵。
\end{answerbox}
\begin{answerbox}[2021--2022 填空 5]
题：向量 $(5,6)^T$ 在基 $(1,2)^T,(3,4)^T$ 下的坐标。

\textbf{【考点】} 基础计算

\textbf{【速解】}

\[x(1,2)^T+y(3,4)^T=(5,6)^T.\]
解得
\textbf{【答案】}

\[[v]_{\alpha}=(-1,2)^T.\]
\end{answerbox}
\begin{answerbox}[2022--2023 解答六：三维基与过渡矩阵]
题：$\alpha_1=(1,2,1)^T,\alpha_2=(1,3,2)^T,\alpha_3=(1,a,3)^T$ 是一组基，$\beta=(1,1,1)^T$ 在该基下坐标为 $(b,c,1)^T$。求 $a,b,c$，并求新旧基过渡矩阵。

\textbf{【题目分析】}

先用
\[\beta=b\alpha_1+c\alpha_2+\alpha_3\]

\textbf{【详细解答】}

比较分量：
\[
(1,1,1)^T=b(1,2,1)^T+c(1,3,2)^T+(1,a,3)^T.
\]
解得
\[
a=3,\qquad b=2,\qquad c=-2.
\]
又
\[
(\alpha_2,\alpha_3,\beta)=(\alpha_1,\alpha_2,\alpha_3)
C,
\]
其中
\[
C=
\begin{pmatrix}
0&0&2\\
1&0&-2\\
0&1&1
\end{pmatrix}.
\]
因为 \(|C|\ne0\)，所以 \(\alpha_2,\alpha_3,\beta\) 是一组基。

\textbf{【最终答案】}

题目要求由 $(\alpha_2,\alpha_3,\beta)$ 到 $(\alpha_1,\alpha_2,\alpha_3)$，故过渡矩阵取 $C^{-1}$。
\[
M=C^{-1}=
\begin{pmatrix}
1&1&0\\
-\frac12&0&1\\
\frac12&0&0
\end{pmatrix}.
\]

\textbf{【方法总结】}

\((\beta)=(\alpha)C\) 时，\(C\) 的第 \(j\) 列是新基第 \(j\) 个向量在旧基下的坐标；由新基到旧基的坐标变换矩阵为 \(C\)，由旧基到新基为 \(C^{-1}\)。
\end{answerbox}
\begin{answerbox}[2023--2024 填空 5]
题：已知 $\gamma=k_1\alpha_1+k_2\alpha_2+k_3\alpha_3$，且 $\gamma^T\alpha_i=\beta^T\alpha_i$，求 $k_1^2+k_2^2+k_3^2$。

\textbf{【考点】} 基础计算

\textbf{【速解】}

因为 $\alpha_1,\alpha_2,\alpha_3$ 是基，条件等价于
\[(\gamma-\beta)^T\alpha_i=0,\ i=1,2,3,\]
$\Rightarrow$ $\gamma=\beta$。求 $\beta$ 在 $\alpha$ 基下坐标，
\textbf{【答案】}

\[k_1^2+k_2^2+k_3^2=9.\]
\end{answerbox}
\begin{answerbox}[2023--2024 解答五：一般 n 维基变换]
题：设向量组（I）：$\alpha_1,\alpha_2,\cdots,\alpha_n$ 是向量空间 $V$ 的一组基。
（1）设 $\beta_1=\alpha_1,\ \beta_i=\alpha_1+\alpha_i\ (2\le i\le n)$，证明向量组（II）：$\beta_1,\beta_2,\cdots,\beta_n$ 也是 $V$ 的一组基；
（2）求基（II）到基（I）的过渡矩阵；
（3）向量 $\alpha$ 在基（I）下的坐标为 $(1,1,\cdots,1)^T$，求 $\alpha$ 在基（II）下的坐标。

\textbf{【题目分析】}

\[
(\beta_1,\ldots,\beta_n)=(\alpha_1,\ldots,\alpha_n)
\begin{pmatrix}
1&1&\cdots&1\\
0&1&\cdots&0\\
\vdots& &\ddots&\vdots\\
0&0&\cdots&1
\end{pmatrix}.
\]

\textbf{【详细解答】}

\[
(\beta_1,\ldots,\beta_n)=(\alpha_1,\ldots,\alpha_n)
\begin{pmatrix}
1&1&\cdots&1\\
0&1&\cdots&0\\
\vdots& &\ddots&\vdots\\
0&0&\cdots&1
\end{pmatrix}.
\]

\textbf{【最终答案】}

该矩阵行列式为 $1$，$\Rightarrow$ $\beta$ 是基。若 $[\alpha]_{\text{基 I}}=(1,1,\ldots,1)^T$，则
\[[\alpha]_{\text{基 II}}=(2-n,1,\ldots,1)^T.\]
\end{answerbox}
\begin{answerbox}[2024--2025 填空 4]
题：设 $\alpha_1,\alpha_2,\alpha_3$ 是三维向量空间 $\mathbb R^3$ 的一个基，则由基
\[
\alpha_1,\ 2\alpha_2,\ 3\alpha_3
\]
到基
\[
\alpha_1+\alpha_2,\ \alpha_2+\alpha_3,\ \alpha_3+\alpha_1
\]
的过渡矩阵为 $\underline{\qquad}$。

\textbf{【考点】} 基础计算

\textbf{【速解】}

以原始 $\alpha$ 为坐标，旧基矩阵
\[D=\operatorname{diag}(1,2,3),\]
新基矩阵
\[C=\begin{pmatrix}1&0&1\\1&1&0\\0&1&1\end{pmatrix}.\]
若采用“旧坐标 = 过渡矩阵 × 新坐标”，则
\textbf{【答案】}

\[
P=D^{-1}C=
\begin{pmatrix}
1&0&1\\
\frac12&\frac12&0\\
0&\frac13&\frac13
\end{pmatrix}.
\]
若按参考答案“由基 \((\alpha_1,2\alpha_2,3\alpha_3)\) 到基 \((\alpha_1+\alpha_2,\alpha_2+\alpha_3,\alpha_3+\alpha_1)\)”的坐标方向，则过渡矩阵写为上式。
\end{answerbox}
\begin{answerbox}[2024--2025 解答五：坐标相同的全部向量]
题：$\beta_1=\alpha_1-\alpha_2,\ \beta_2=2\alpha_2-\alpha_3,\ \beta_3=3\alpha_3$。证明 $\beta$ 为基，并求在两组基下坐标相同的全部向量。

\textbf{【题目分析】}

\[
(\beta_1,\beta_2,\beta_3)=(\alpha_1,\alpha_2,\alpha_3)
\begin{pmatrix}
1&0&0\\
-1&2&0\\
0&-1&3
\end{pmatrix},
\]

\textbf{【详细解答】}

行列式 $=6\ne0$，$\Rightarrow$为基。设共同坐标 $x$，则
\[(\alpha)x=(\beta)x=(\alpha)Cx,\]

因此
\[(C-E)x=0.\]

\textbf{【最终答案】}

解出 $x$ 后，向量为
\[\gamma=(\alpha_1,\alpha_2,\alpha_3)x.\]
\end{answerbox}
\bigskip\hrule\bigskip

\section{特征值、特征向量、相似对角化、实对称矩阵}

\begin{identifybox}[高亮公式]
$A\alpha=\lambda\alpha\Rightarrow f(A)\alpha=f(\lambda)\alpha$。\\
实对称矩阵：不同特征值对应特征向量正交；一定可正交对角化。\\
对角化步骤：求特征值 -> 求特征向量 -> 组 $P$ 或正交单位化组 $Q$。\\
\end{identifybox}

\begin{templatebox}[考场模板]

\begin{enumerate}
\item 看到“实对称”先写：可正交对角化。
\item 已知特征向量时，用 $A\alpha_i=\lambda_i\alpha_i$ 反求未知参数。
\item 正交矩阵 $Q$ 的列必须是单位正交特征向量。
\item 求 $A$：写 $A=Q\Lambda Q^T$ 或 $A=P\Lambda P^{-1}$。
\end{enumerate}

\end{templatebox}
\subsection{真题融合}

\begin{answerbox}[2020--2021 填空 5]
题：设 $A$ 为 $3$ 阶实对称矩阵，
\[
\alpha_1=\begin{pmatrix}2\\a+4\\-3\end{pmatrix},\quad
\alpha_2=\begin{pmatrix}a+1\\1\\1\end{pmatrix},
\]
且 $A\alpha_1=\alpha_1,\ A\alpha_2=2\alpha_2$，则 $a=\underline{\qquad}$。

\textbf{【考点】} 基础计算

\textbf{【速解】}

实对称矩阵对应不同特征值的特征向量正交。直接令
\textbf{【答案】}

\[\alpha_1^T\alpha_2=0\]
解参数， $-2$。
\end{answerbox}
\begin{answerbox}[2020--2021 解答五：实对称矩阵补特征向量]
题：三阶实对称矩阵 $A$ 的三个特征值为 $\lambda_1=\lambda_2=1,\ \lambda_3=2$，且对应于 $\lambda_1=\lambda_2=1$ 的特征向量为
\[
\alpha_1=(1,1,-1)^T,\quad \alpha_2=(-2,-3,3)^T.
\]
（1）求 $A$ 的对应于特征值 $\lambda_3=2$ 的特征向量；
（2）求正交矩阵 $Q$ 及对角形矩阵 $\Lambda$，使 $Q^{-1}AQ=\Lambda$。

\textbf{【题目分析】}

$\lambda=2$ 的特征向量必须与已知两个向量都正交。设 $\xi=(x_1,x_2,x_3)^T$，解
\[\xi^T\xi_1=0,\quad \xi^T\xi_2=0.\]

\textbf{【详细解答】}

$\lambda=2$ 的特征向量必须与已知两个向量都正交。设 $\xi=(x_1,x_2,x_3)^T$，解
\[\xi^T\xi_1=0,\quad \xi^T\xi_2=0.\]

\textbf{【最终答案】}

得 $\xi=k(0,1,1)^T$。把三个特征向量正交单位化，作
\[Q=(e_1,e_2,e_3),\quad Q^TAQ=\operatorname{diag}(-1,-1,2).\]
\end{answerbox}
\begin{answerbox}[2022--2023 填空 5]
题：三阶矩阵满足 $|A-E|=|A-2E|=|A-3E|=0$，求 $|A-4E|$。

\textbf{【考点】} 基础计算

\textbf{【速解】}

特征值为 $1,2,3$。
\textbf{【答案】}

\[|A-4E|=(1-4)(2-4)(3-4)=-6.\]
\end{answerbox}
\begin{answerbox}[2022--2023 解答四：四阶矩阵对角化]
题：设四阶矩阵
\[
A=\begin{pmatrix}
3&2&2&2\\
a&1&b&c\\
d&e&1&f\\
g&h&k&1
\end{pmatrix}
\]
对应特征值 $\lambda=3$ 有三个线性无关的特征向量，其中 $a,b,c,d,e,f,g,h,k$ 为常数。
（1）求矩阵 $A$；
（2）求可逆矩阵 $P$ 和对角形矩阵 $\Lambda$，使 $P^{-1}AP=\Lambda$。

\textbf{【题目分析】}

先由 $(A-3E)x=0$ 的维数为 $3$ 写出参数限制，求出矩阵 $A$。补另一个特征值及特征向量。

\textbf{【详细解答】}

由 \(r(3E-A)=1\) 得
\[
A=
\begin{pmatrix}
3&2&2&2\\
0&1&-2&-2\\
0&-2&1&-2\\
0&-2&-2&1
\end{pmatrix}.
\]
由 \(\operatorname{tr}(A)=6\)，第四个特征值为 \(-3\)。可取
\[
\xi_1=(1,0,0,0)^T,\quad
\xi_2=(0,-1,1,0)^T,\quad
\xi_3=(0,-1,0,1)^T,
\]
\[
\xi_4=(-1,1,1,1)^T.
\]

\textbf{【最终答案】}

\[
P=
\begin{pmatrix}
1&0&0&-1\\
0&-1&-1&1\\
0&1&0&1\\
0&0&1&1
\end{pmatrix},
\qquad
\Lambda=\operatorname{diag}(3,3,3,-3).
\]
\[
P^{-1}AP=\Lambda.
\]

\textbf{【方法总结】}

\(P\) 的列向量顺序必须与 \(\Lambda\) 的对角元顺序一致。
\end{answerbox}
\begin{answerbox}[2023--2024 填空 6]
题：设二阶矩阵 $A$ 有两个不同特征值，$\alpha_1,\alpha_2$ 是 $A$ 的线性无关的特征向量，且满足
\[A^2(\alpha_1+\alpha_2)=\alpha_1+\alpha_2.\]
求 $|A|$。

\textbf{【考点】} 基础计算

\textbf{【速解】}

设对应特征值为 $\lambda_1,\lambda_2$。由线性无关得
\[\lambda_1^2=1,\quad \lambda_2^2=1.\]
又两个特征值不同，得特征值为 $1,-1$。
\textbf{【答案】}

\[|A|=-1.\]
\end{answerbox}
\begin{answerbox}[2023--2024 解答六：实对称矩阵由多项式确定]
题：三阶实对称矩阵 $A$ 可逆，满足 $A^3-3A^2+2A=O$，且 $Ax=x$ 的基础解系为 $(1,1,0)^T$。求全部特征值、正交矩阵 $Q$ 与 $A$。

\textbf{【题目分析】}

\[A(A-E)(A-2E)=O.\]

\textbf{【详细解答】}

因 $A$ 可逆，特征值只可能为 $1,2$。又 $\lambda=1$ 的特征子空间维数为 $1$，$\Rightarrow$特征值为
\[1,2,2.\]

取
\[e_1=\frac1{\sqrt2}(1,1,0)^T,\]

\textbf{【最终答案】}

取
\[
\eta_1=\frac1{\sqrt2}(1,1,0)^T,\qquad
\eta_2=\frac1{\sqrt2}(-1,1,0)^T,\qquad
\eta_3=(0,0,1)^T.
\]
令
\[
Q=(\eta_1,\eta_2,\eta_3)=
\begin{pmatrix}
\frac1{\sqrt2}&-\frac1{\sqrt2}&0\\
\frac1{\sqrt2}&\frac1{\sqrt2}&0\\
0&0&1
\end{pmatrix},
\quad
\Lambda=\operatorname{diag}(1,2,2).
\]
\[
Q^TAQ=\Lambda,\qquad
A=Q\Lambda Q^T=
\begin{pmatrix}
\frac32&-\frac12&0\\
-\frac12&\frac32&0\\
0&0&2
\end{pmatrix}.
\]

\textbf{【方法总结】}

实对称矩阵正交对角化时，\(Q\) 的列必须是单位正交特征向量。
\end{answerbox}
\begin{answerbox}[2024--2025 填空 3]
题：设三阶实对称矩阵 $A$ 的行列式 $|A|=6$，并且
\[
|A+E|+|A-3E|=0.
\]
求二次型
\[
f(x_1,x_2,x_3)=X^TAX,\qquad X=(x_1,x_2,x_3)^T
\]
的正惯性指数，这里 $E$ 为 $3$ 阶单位矩阵。

\textbf{【考点】} 基础计算

\textbf{【速解】}

特征值满足
\[\lambda^3-7\lambda-6=0=(\lambda-3)(\lambda+1)(\lambda+2).\]
乘积 $3\cdot(-1)\cdot(-2)=6$ 与 $|A|$ 符合，正特征值只有一个，$\Rightarrow$正惯性指数为
\textbf{【答案】}

\[1.\]
\end{answerbox}
\begin{answerbox}[2024--2025 填空 6]
题：$A$ 有特征值 $-1,1,2,\ldots,2$，$\xi_1,\xi_2$ 分别属于 $-1,1$，问 $7\xi_1-3\xi_2$ 是 $A^{100}$ 的哪个特征值的特征向量。

\textbf{【考点】} 基础计算

\textbf{【速解】}

\[A^{100}\xi_1=(-1)^{100}\xi_1=\xi_1,\quad A^{100}\xi_2=\xi_2.\]
所以
\textbf{【答案】}

\[A^{100}(7\xi_1-3\xi_2)=1\cdot(7\xi_1-3\xi_2),\]
 $1$。
\end{answerbox}
\begin{answerbox}[2024--2025 解答六：由特征向量和第一行求实对称矩阵]
题：三阶实方阵 $A$ 第一行为 $(1,2,0)$，且
\[\alpha_1=(-1,0,1)^T,\quad \alpha_2=(1,1,1)^T,\quad \alpha_3=(1,-2,1)^T\]
为三个特征向量。求正交矩阵 $Q$ 与 $A$。

\textbf{【题目分析】}

先由第一行计算特征值：
\[\lambda_i=\frac{(1,2,0)\alpha_i}{(\alpha_i)_1},\]

\textbf{【详细解答】}

先由第一行计算特征值：
\[\lambda_i=\frac{(1,2,0)\alpha_i}{(\alpha_i)_1},\]
得
\[
\lambda_1=1,\qquad \lambda_2=3,\qquad \lambda_3=-3.
\]
将三个特征向量单位化：
\[
\eta_1=\frac1{\sqrt2}(-1,0,1)^T,\quad
\eta_2=\frac1{\sqrt3}(1,1,1)^T,\quad
\eta_3=\frac1{\sqrt6}(1,-2,1)^T.
\]

\textbf{【最终答案】}

\[
Q=
\begin{pmatrix}
-\frac1{\sqrt2}&\frac1{\sqrt3}&\frac1{\sqrt6}\\
0&\frac1{\sqrt3}&-\frac2{\sqrt6}\\
\frac1{\sqrt2}&\frac1{\sqrt3}&\frac1{\sqrt6}
\end{pmatrix},
\qquad
\Lambda=\operatorname{diag}(1,3,-3).
\]
\[
Q^TAQ=\Lambda.
\]
按参考答案计算
\[
A=Q\Lambda Q^T=
\begin{pmatrix}
1&2&0\\
2&-1&2\\
0&2&1
\end{pmatrix}.
\]

\textbf{【方法总结】}

已知第一行和特征向量时，先用第一行点乘特征向量求特征值，再用 \(A=Q\Lambda Q^T\) 求矩阵。
\end{answerbox}
\bigskip\hrule\bigskip

\section{二次型：正定、惯性指数、标准形、正交变换}

\begin{identifybox}[高亮公式]
二次型 $f=x^TAx$ 必须先写对称矩阵 $A$。\\
正定：顺序主子式全正；负定：$(-1)^kD_k>0$。\\
实对称正交化：求特征值和单位特征向量，$x=Qy$，标准形为 $\lambda_1y_1^2+\cdots+\lambda_ny_n^2$。\\
\end{identifybox}

\begin{templatebox}[考场模板]

\begin{enumerate}
\item 把交叉项系数除以 $2$ 写入矩阵。
\item 问正定/负定：优先顺序主子式。
\item 问正交变换：求特征值、特征向量、单位化。
\item 问是否存在正交变换把 $f$ 化为 $g$：比较特征值；只合同则比较正负惯性指数。
\end{enumerate}

\end{templatebox}
\subsection{真题融合}

\begin{answerbox}[2020--2021 填空 6]
题：二次型
\[
f(x_1,x_2,x_3)=
\begin{pmatrix}x_1&x_2&x_3\end{pmatrix}
\begin{pmatrix}
1&2&2\\
0&1&2\\
0&0&1
\end{pmatrix}
\begin{pmatrix}x_1\\x_2\\x_3\end{pmatrix}
\]
的秩为 $\underline{\qquad}$。

\textbf{【考点】} 基础计算

\textbf{【速解】}

二次型秩就是对应对称矩阵的秩。行化简得
\textbf{【答案】}

\[r(A)=1.\]
\end{answerbox}
\begin{answerbox}[2020--2021 解答七：循环平方和二次型正定]
题：
\[f=(x_1+a_1x_2)^2+(x_2+a_2x_3)^2+\cdots+(x_n+a_nx_1)^2.\]
证明半正定，并求正定条件。

\textbf{【题目分析】}

由平方和结构
\[f\ge0.\]

\textbf{【详细解答】}

正定要求齐次方程组
\[x_1+a_1x_2=0,\ldots,x_n+a_nx_1=0\]

\textbf{【最终答案】}

只有零解，即系数行列式非零。条件：
\[a_1a_2\cdots a_n\ne(-1)^n.\]
\end{answerbox}
\begin{answerbox}[2021--2022 填空 6]
题：
\[f=(x_1+ax_2)^2+(x_2+bx_3)^2+(x_3+cx_1)^2\]
何时正定。

\textbf{【考点】} 基础计算

\textbf{【速解】}

同上，正定等价于线性变换可逆：
\textbf{【答案】}

\[abc\ne -1.\]
\end{answerbox}
\begin{answerbox}[2021--2022 解答五：三元二次型正交标准形]
题：
\[f=ax_1^2+2x_2^2-2x_3^2+2bx_1x_3,\quad a,b>0,\]
已知矩阵特征值之和为 $1$、积为 $-12$，求 $a,b$ 并正交化。

\textbf{【题目分析】}

先写矩阵
\[A=\begin{pmatrix}a&0&b\\0&2&0\\b&0&-2\end{pmatrix}.\]

\textbf{【详细解答】}

由迹与行列式：
\[a=1,\quad 2(-2a-b^2)=-12.\]

因 $b>0$，得
\[a=1,\quad b=2.\]

\textbf{【最终答案】}

此时特征值为 $2,2,-3$，正交变换后标准形为
\[f=2y_1^2+2y_2^2-3y_3^2.\]
可取正交矩阵
\[
Q=
\begin{pmatrix}
\frac2{\sqrt5}&0&\frac1{\sqrt5}\\
0&1&0\\
\frac1{\sqrt5}&0&-\frac2{\sqrt5}
\end{pmatrix},
\qquad x=Qy.
\]

\textbf{【方法总结】}

正交化题必须写对称矩阵、特征值、单位特征向量和标准形。
\end{answerbox}
\begin{answerbox}[2022--2023 填空 6]
题：
\[f=tx_1^2+tx_2^2+(t-2)x_3^2+2x_1x_2\]
负定，求 $t$ 范围。

\textbf{【考点】} 基础计算

\textbf{【速解】}

写对称矩阵
\[A=\begin{pmatrix}t&1&0\\1&t&0\\0&0&t-2\end{pmatrix}.\]
负定条件：
\[D_1<0,\quad D_2>0,\quad D_3<0.\]
即
\[t<0,\quad t^2-1>0,\quad (t^2-1)(t-2)<0.\]
合并得
\textbf{【答案】}

\[t<-1.\]
\end{answerbox}
\begin{answerbox}[2022--2023 解答七：实对称矩阵与二次型]
题：设三阶实对称矩阵 $A$ 的特征值为 $1,1,4$，向量
\[
\xi_1=(-1,1,0)^T,\quad \xi_2=(-1,0,1)^T
\]
是 $A$ 的对应特征值 $\lambda=1$ 的特征向量。记 $X=(x_1,x_2,x_3)^T$，二次型
\[
f(x_1,x_2,x_3)=X^TAX.
\]
（1）判断二次型 $f(x_1,x_2,x_3)$ 是否正定并说明理由；
（2）用正交变换将二次型 $f(x_1,x_2,x_3)$ 化为标准型，写出所用的正交变换；
（3）求出二次型 $f(x_1,x_2,x_3)$。

\textbf{【题目分析】}

特征值全正，$\Rightarrow$正定。把两个已知特征向量正交单位化，取它们的正交补作 $\lambda=4$ 的单位特征向量，得

\textbf{【详细解答】}

参考答案将 \(\lambda=1\) 的两个特征向量正交化、单位化：
\[
\eta_1=\left(-\frac1{\sqrt2},\frac1{\sqrt2},0\right)^T,\qquad
\eta_2=\left(-\frac1{\sqrt6},-\frac1{\sqrt6},\frac2{\sqrt6}\right)^T.
\]
\(\lambda=4\) 的单位特征向量取
\[
\eta_3=\left(\frac1{\sqrt3},\frac1{\sqrt3},\frac1{\sqrt3}\right)^T.
\]
令 \(Q=(\eta_1,\eta_2,\eta_3)\)。

\textbf{【最终答案】}

\[
Q=
\begin{pmatrix}
-\frac1{\sqrt2}&-\frac1{\sqrt6}&\frac1{\sqrt3}\\
\frac1{\sqrt2}&-\frac1{\sqrt6}&\frac1{\sqrt3}\\
0&\frac2{\sqrt6}&\frac1{\sqrt3}
\end{pmatrix},
\qquad
Q^TAQ=\operatorname{diag}(1,1,4).
\]
在正交变换 \(X=QY\) 下
\[
f=y_1^2+y_2^2+4y_3^2.
\]
并且
\[
A=Q\operatorname{diag}(1,1,4)Q^T=
\begin{pmatrix}
2&1&1\\
1&2&1\\
1&1&2
\end{pmatrix},
\]
\[
f(x_1,x_2,x_3)=2x_1^2+2x_2^2+2x_3^2+2x_1x_2+2x_1x_3+2x_2x_3.
\]

\textbf{【方法总结】}

同一特征值下给出的特征向量若不正交，先施密特正交化，再单位化。
\end{answerbox}
\begin{answerbox}[2023--2024 解答七：合同与正交的区别]
题：
\[f=x_1^2+2x_2^2+2x_3^2+2x_1x_2-2x_1x_3,\]
\[g=y_1^2+y_2^2+y_3^2+2y_2y_3.\]
求 $f$ 规范形，找可逆变换使 $f$ 化为 $g$，并判断是否存在正交变换。

\textbf{【题目分析】}

二次型矩阵分别为
\[
A=\begin{pmatrix}1&1&-1\\1&2&0\\-1&0&2\end{pmatrix},
\qquad
B=\begin{pmatrix}1&0&0\\0&1&1\\0&1&1\end{pmatrix}.
\]
参考答案给出
\[
f=(x_1+x_2-x_3)^2+(x_2+x_3)^2,
\]
故 \(f\) 的规范形为
\[
z_1^2+z_2^2.
\]

\textbf{【详细解答】}

令
\[
y_1=x_1+x_2-x_3,\qquad y_2=x_2,\qquad y_3=x_3,
\]
即
\[
x=Py,\qquad
P=
\begin{pmatrix}
1&-1&1\\
0&1&0\\
0&0&1
\end{pmatrix}.
\]
则
\[
P^TAP=B,
\]
故可逆线性变换 \(x=Py\) 可将 \(f\) 化为 \(g\)。

\textbf{【最终答案】}

\[
f\text{ 的规范形为 }z_1^2+z_2^2.
\]
\[
P=
\begin{pmatrix}
1&-1&1\\
0&1&0\\
0&0&1
\end{pmatrix}
\]
可使 \(x=Py\) 将 \(f\) 化为 \(g\)。又
\[
\operatorname{tr}(A)=5,\qquad \operatorname{tr}(B)=3,
\]
故 \(A,B\) 不相似，不存在正交变换把 \(f\) 化为 \(g\)。

\textbf{【方法总结】}

合同只保惯性指数；正交变换要求矩阵正交相似，必须比较特征值或迹等相似不变量。
\end{answerbox}
\begin{answerbox}[2024--2025 解答七：秩为 2 的二次型]
题：
\[f=(1-a)x_1^2+(1-a)x_2^2+2x_3^2+2(1+a)x_1x_2\]
秩为 $2$。求 $a$，并正交化。

\textbf{【题目分析】}

矩阵为
\[A=\begin{pmatrix}1-a&1+a&0\\1+a&1-a&0\\0&0&2\end{pmatrix}.\]

\textbf{【详细解答】}

秩为 $2$ 等价于一个特征值为 $0$，得
\[a=0.\]

此时特征值为 $2,2,0$。取单位特征向量组成
\[
P=\begin{pmatrix}
\frac{\sqrt2}{2}&0&\frac{\sqrt2}{2}\\
\frac{\sqrt2}{2}&0&-\frac{\sqrt2}{2}\\
0&1&0
\end{pmatrix},
\]

\textbf{【最终答案】}

标准形为
\[f=2y_1^2+2y_2^2.\]
规范形为
\[
z_1^2+z_2^2.
\]
正惯性指数为 \(2\)，负惯性指数为 \(0\)，秩为 \(2\)。

\textbf{【方法总结】}

秩为 \(2\) 表示只有两个非零平方项；正交标准形保留特征值，规范形只保留正负号。
\end{answerbox}
\bigskip\hrule\bigskip

\section{证明题：固定套路}

\begin{identifybox}[高亮公式]
证明“仍为基础解系”：先证明都是解，证明线性无关，个数等于维数。\\
证明“可逆”：常用 $|A|\ne0$、零空间只有 $0$、特征值不含 $0$。\\
证明“可对角化”：凑够 $n$ 个线性无关特征向量。\\
\end{identifybox}

\subsection{真题融合}

\begin{answerbox}[2020--2021 证明八(2)：\ensuremath{A^T=A^*} 推正交矩阵]
题：$A$ 为 $n$ 阶非零实矩阵，$n>2$，$A^T=A^*$，证明 $A$ 为正交矩阵。

\textbf{【证明思路】}

\[AA^T=AA^*=|A|E.\]

\textbf{【详细证明】}

先由非零行说明 $|A|>0$。两边取行列式：
\[|A|^2=|A|^n.\]

因 $n>2$ 且 $|A|>0$，得 $|A|=1$，故
\[AA^T=E,\]

\textbf{【最终答案】}

$\Rightarrow$ $A$ 正交。
\end{answerbox}
\begin{answerbox}[2021--2022 证明一：由递推关系证明矩阵可逆]
题：$\alpha,A\alpha,A^2\alpha$ 线性无关，且 $A^3\alpha=3A\alpha-2A^2\alpha$，证明由若干列构成的矩阵 $B$ 可逆。

\textbf{【证明思路】}

把 $B$ 的列全部写成 $\alpha,A\alpha,A^2\alpha$ 的线性组合：
\[B=(\alpha,A\alpha,A^2\alpha)C.\]

\textbf{【详细证明】}

左边第一矩阵可逆，算
\[|C|\ne0.\]

\textbf{【最终答案】}

因此 $B$ 可逆。
\end{answerbox}
\begin{answerbox}[2021--2022 证明二：幂等矩阵相似标准形]
题：$A^2=A,\ r(A)=r$，证明
\[A\sim \begin{pmatrix}E_r&0\\0&0\end{pmatrix}.\]

\textbf{【证明思路】}

$A^2=A$ 的特征值只可能为 $0,1$。又 $r(A)=r$，说明 $\lambda=1$ 的特征向量取 $r$ 个，$\lambda=0$ 的特征向量取 $n-r$ 个。合成 $n$ 个线性无关特征向量，对角化为上式。

\textbf{【详细证明】}

$A^2=A$ 的特征值只可能为 $0,1$。又 $r(A)=r$，说明 $\lambda=1$ 的特征向量取 $r$ 个，$\lambda=0$ 的特征向量取 $n-r$ 个。合成 $n$ 个线性无关特征向量，对角化为上式。

\textbf{【最终答案】}

$A^2=A$ 的特征值只可能为 $0,1$。又 $r(A)=r$，说明 $\lambda=1$ 的特征向量取 $r$ 个，$\lambda=0$ 的特征向量取 $n-r$ 个。合成 $n$ 个线性无关特征向量，对角化为上式。
\end{answerbox}
\begin{answerbox}[2022--2023 证明一：伴随矩阵与基础解系]
题：四阶矩阵 $A=(\alpha_1,\alpha_2,\alpha_3,\alpha_4)$，$AX=0$ 的基础解系为 $(1,2,0,0)^T$，证明若干列向量是 $A^*x=0$ 的基础解系。

\textbf{【证明思路】}

由基础解系可知
\[\alpha_1+2\alpha_2=0,\quad r(A)=3.\]

\textbf{【详细证明】}

由基础解系可知
\[\alpha_1+2\alpha_2=0,\quad r(A)=3.\]

\textbf{【最终答案】}

$\Rightarrow$ $r(A^*)=1$，$A^*x=0$ 的解空间维数为 $3$。验证三个向量都是解且线性无关，即为基础解系。
\end{answerbox}
\begin{answerbox}[2022--2023 证明二：二次多项式条件下可对角化]
题：$A^2-(a+b)A+abE=O,\ a\ne b$。证明
\[r(A-aE)+r(A-bE)=n,\]
且 $A$ 可对角化。

\textbf{【证明思路】}

\[(A-aE)(A-bE)=O.\]

\textbf{【详细证明】}

利用
\[\mathbb R^n=N(A-aE)\oplus N(A-bE)\]

\textbf{【最终答案】}

证明两个特征子空间维数之和为 $n$，$\Rightarrow$有 $n$ 个线性无关特征向量，$A$ 可对角化。
\end{answerbox}
\begin{answerbox}[2023--2024 证明一：非齐次解组线性无关]
题：$b\ne0$，$\eta_0$ 是 $Ax=b$ 的一个解，$\xi_1,\ldots,\xi_{n-r}$ 是 $Ax=0$ 的基础解系，令 $\eta_i=\eta_0+\xi_i$。证明 $\eta_0,\eta_1,\ldots,\eta_{n-r}$ 线性无关。

\textbf{【证明思路】}

设
\[k_0\eta_0+k_1\eta_1+\cdots+k_s\eta_s=0.\]

\textbf{【详细证明】}

代入 $\eta_i=\eta_0+\xi_i$：
\[(k_0+\cdots+k_s)\eta_0+k_1\xi_1+\cdots+k_s\xi_s=0.\]

两边左乘 $A$：
\[(k_0+\cdots+k_s)b=0.\]

\textbf{【最终答案】}

因 $b\ne0$，得 $k_0+\cdots+k_s=0$。由 $\xi_i$ 线性无关得 $k_1=\cdots=k_s=0$，$\Rightarrow$ $k_0=0$。
\end{answerbox}
\begin{answerbox}[2023--2024 证明二：\ensuremath{CC^T=O} 与迹技巧]
题：$A,B,C$ 为 $n$ 阶实矩阵。证明：若 $CC^T=O$，则 $C=O$；若 $AB-BA=CC^T$，则 $AB=BA$。

\textbf{【证明思路】}

\begin{enumerate}
\item $CC^T=O$ 时，第 $i$ 个对角元为 $C$ 第 $i$ 行元素平方和，$\Rightarrow$每行全零，$C=O$。
\item 两边取迹：
\end{enumerate}
\[\operatorname{tr}(AB-BA)=0.\]

\textbf{【详细证明】}

故
\[\operatorname{tr}(CC^T)=0.\]

\textbf{【最终答案】}

平方和为 $0$，得 $C=O$，$\Rightarrow$ $AB-BA=O$。
\end{answerbox}
\begin{answerbox}[2024--2025 证明一：Vandermonde 型线性无关]
题：$A$ 有三个不同特征值 $\lambda_1,\lambda_2,\lambda_3$，对应特征向量 $\alpha_1,\alpha_2,\alpha_3$，$\beta=\alpha_1+\alpha_2+\alpha_3$。证明 $\beta,A\beta,A^2\beta$ 线性无关。

\textbf{【证明思路】}

设
\[k_1\beta+k_2A\beta+k_3A^2\beta=0.\]

\textbf{【详细证明】}

展开：
\[\sum_{i=1}^3(k_1+k_2\lambda_i+k_3\lambda_i^2)\alpha_i=0.\]

不同特征值对应特征向量线性无关，得
\[k_1+k_2\lambda_i+k_3\lambda_i^2=0,\quad i=1,2,3.\]

系数矩阵是 Vandermonde 矩阵，行列式
\[\prod_{i<j}(\lambda_j-\lambda_i)\ne0,\]

\textbf{【最终答案】}

$\Rightarrow$ $k_1=k_2=k_3=0$。
\end{answerbox}
\begin{answerbox}[2024--2025 证明二：秩一矩阵逆公式]
题：$r(A)=1,\ \operatorname{tr}(A)=2$，证明 $I+A$ 可逆且
\[(I+A)^{-1}=I-\frac13A.\]

\textbf{【证明思路】}

秩为 $1$，非零特征值只有一个；迹为 $2$，$\Rightarrow$特征值为 $2,0,\ldots,0$。$\Rightarrow$ $I+A$ 的特征值为 $3,1,\ldots,1$，可逆。由秩一矩阵满足
\[A^2=\operatorname{tr}(A)A=2A,\]

\textbf{【详细证明】}

秩为 $1$，非零特征值只有一个；迹为 $2$，$\Rightarrow$特征值为 $2,0,\ldots,0$。$\Rightarrow$ $I+A$ 的特征值为 $3,1,\ldots,1$，可逆。由秩一矩阵满足
\[A^2=\operatorname{tr}(A)A=2A,\]

\textbf{【最终答案】}

直接验算：
\[
(I+A)\left(I-\frac13A\right)
=I+\frac23A-\frac13A^2
=I.
\]
\end{answerbox}
\bigskip\hrule\bigskip

\section{考前速背清单}
\begin{scorebox}[考前速背清单]

\begin{enumerate}
\item 行列式题：先看倍乘、伴随、列变换，不硬展开。
\item 秩题：先写矩阵，行化简，主元列对应极大无关组。
\item 方程组题：永远先比较 $r(A)$ 与 $r(A,b)$。
\item 通解题：必须写“特解 + 基础解系”。
\item 基变换题：先写 $(\beta)=(\alpha)C$，决定取 $C$ 还是 $C^{-1}$。
\item 实对称题：正交对角化；不同特征值的特征向量正交。
\item 二次型题：交叉项除以 $2$ 写矩阵；正交标准形看特征值。
\item 证明题：先把目标翻译成“线性无关、可逆、秩、特征值、迹”之一。
\end{enumerate}
\end{scorebox}

\end{document}



